\section{Tournaments and a vertex of maximum outdegree}\label{ch07:sec:tournament}
In a round-robin competition, every pair of players meets exactly once, and every game has a winner. We can record the results as an oriented graph: the players are the vertices, and $u\to v$ means that $u$ beat $v$.

\par\noindent\begin{minipage}{\linewidth}
\begin{definition}[Tournament]\label{ch07:def:tournament}
A \emph{tournament} is an oriented graph in which every pair of different vertices is joined by exactly one directed edge: for all vertices $u\ne v$, exactly one of $u\to v$ and $v\to u$ holds.
\end{definition}
\end{minipage}\par

Here is a tournament with four players $p$, $q$, $r$, $s$ and the six results $p\to q$, $p\to r$, $q\to r$, $q\to s$, $r\to s$, and $s\to p$, one for each pair:
\begin{center}
\begin{tikzpicture}[>=Stealth,every node/.style={circle,draw=ink,minimum size=6mm,inner sep=1pt},line width=.7pt]
\node(p)at(0,1.8){$p$};\node(q)at(2.4,1.8){$q$};\node(r)at(2.4,0){$r$};\node(s)at(0,0){$s$};
\draw[->](p)--(q);\draw[->](p)--(r);\draw[->](q)--(r);\draw[->](q)--(s);\draw[->](r)--(s);\draw[->](s)--(p);
\end{tikzpicture}
\end{center}
The outdegrees, the numbers of games won, are $2$ for $p$ and $q$, and $1$ for $r$ and $s$.

Call a vertex $v$ a \emph{king} if it can reach every other vertex by a directed path of length at most two: for every other player $u$, either $v$ beat $u$, or $v$ beat someone who beat $u$. In the example, $p$ is a king: it beat $q$ and $r$, and it reaches $s$ through $p\to q\to s$. So is $q$, which beat $r$ and $s$ and reaches $p$ through $q\to s\to p$. The player $r$ is not a king: its only win was against $s$, and $s$ beat only $p$, so within two steps $r$ does not reach $q$. Must every nonempty finite tournament have a king? An extremal choice answers this question.

\par\noindent\begin{minipage}{\linewidth}
\begin{proposition}\label{thm:king}
In a nonempty finite tournament, every vertex of maximum outdegree is a king: it can reach every other vertex by a directed path of length at most two.
\end{proposition}
\end{minipage}\par

A vertex of maximum outdegree always exists, because the outdegrees form a nonempty finite set of natural numbers, which has a largest member. So the proposition shows that every nonempty finite tournament has a king.

\begin{proof}
Let $v$ be any vertex of maximum outdegree, and let $A$ be the set of its out-neighbors, so that the outdegree of $v$ is $|A|$. Let $u$ be any vertex other than $v$. There are two ways in which $u$ might be reached quickly.
\begin{itemize}
\item If $v\to u$, then $u$ is reached in one step.
\item If some $a\in A$ satisfies $a\to u$, then $u$ is reached in two steps by $v\to a\to u$. This is a directed path, because its three vertices are different: $a\ne v$ because $v\to a$ and there are no loops, $u\ne v$ by the choice of $u$, and $a\ne u$ because $a\to u$.
\end{itemize}
Suppose, for contradiction, that neither happens. Since $v\to u$ fails and $u\ne v$, the tournament condition gives $u\to v$. Also $u\notin A$, because $v\to u$ fails, so $u$ is different from every $a\in A$. For each $a\in A$, the tournament condition gives $a\to u$ or $u\to a$, and $a\to u$ has been excluded, so $u\to a$. Thus $u$ beat every member of $A$, and it also beat $v$, which is not in $A$ because a vertex is never its own out-neighbor. So $u$ has at least $|A|+1$ out-neighbors, while $v$ has exactly $|A|$. This contradicts the choice of $v$ as a vertex of maximum outdegree. Therefore one of the two ways always occurs, and $v$ is a king.
\end{proof}

\pauseq{In the example tournament, $s$ won only one game, so it does not have maximum outdegree. Is $s$ a king? What does your answer say about the converse of Proposition~\ref{thm:king}?}{Yes: $s\to p$, $s\to p\to q$, and $s\to p\to r$ reach all three other players within two steps. So the converse, ``every king has maximum outdegree,'' is false, with $s$ as a counterexample: the proposition gives a sufficient condition for being a king, not a necessary one. The kings of the example are $p$, $q$, and $s$.}
